% Part: first-order-logic
% Chapter: proof-systems
% Section: natural-deduction

\documentclass[../../../include/open-logic-section]{subfiles}

\begin{document}

\iftag{FOL}
      {\olfileid[tr]{fol}{prf}{ntd}}
      {\olfileid[tr]{pl}{prf}{ntd}}

\olsection{Doğal Tümdengelim}

Doğal tümdengelim, fiilî akıl yürütmeyi (özellikle matematikçilerin kullandığı
düzenli akıl yürütme biçimini) yansıtmayı amaçlayan !!a{derivation}
sistemidir. Fiilî akıl yürütme birtakım “doğal” örüntülere göre ilerler.
Örneğin durumlara ayırarak ispat, “veya” bağlaçlı bir öncülden hareketle
sonucun veya bileşenlerinin her birinden çıktığını ayrı ayrı göstererek bir
sonuç ortaya koymamızı sağlar. Dolaylı ispat, bir sonucun değillemesinin çelişkiye yol açtığını
göstererek o sonucu ortaya koymamızı sağlar. Koşullu ispat ise artbileşenin
önbileşenden çıktığını göstererek “eğer \dots ise \dots” biçimindeki koşullu
bir iddiayı ortaya koyar. Doğal tümdengelim, bu doğal çıkarımların bazılarının
biçimselleştirilmesidir. Mantıksal bağlaçların
\iftag{FOL}{ ve niceleyicilerin}{} her birinin,
biri giriş diğeri giderme olmak üzere iki kuralı vardır ve bu kuralların her
biri böyle bir doğal çıkarım örüntüsüne karşılık gelir. Örneğin
$\Intro{\lif}$ koşullu ispata, $\Elim{\lor}$ ise durumlara ayırarak ispata
karşılık gelir. Özellikle basit bir kural olan $\Elim{\land}$, $!A \land !B$
ifadesinden $!A$'nın (veya $!B$'nin) çıkarılmasına izin verir.

Doğal tümdengelimi diğer !!{derivation} sistemlerinden ayıran özelliklerden
biri varsayımları kullanmasıdır. Doğal tümdengelimde !!^a{derivation},
!!{formula}s içeren bir ağaçtır. Ağacın kökünde tek bir !!{formula}, ağacın
“yapraklarında” ise sonucun türetildiği !!{formula}s bulunur. Doğal
tümdengelimde yapraktaki bazı !!{formula}s, başka !!{formula}s ile birlikte
!!{derivation} içinde rol oynar ama !!{derivation} sonuca ulaştığında artık
kullanılmaz. Bu, fiilî akıl
yürütmede yalnızca kısa süre geçerli olan hipotezler öne sürme pratiğine
karşılık gelir. Örneğin durumlara ayırarak ispatta veya bileşenlerinin her
birini ayrı bir durumda doğru varsayarız; koşullu ispatta önbileşeni, dolaylı
ispatta ise sonucun değillemesini doğru varsayarız. Hipotetik varsayımları öne
sürüp bir ara adımı elde ettikten sonra ortadan kaldırmak, doğal tümdengelimin
ayırt edici özelliğidir. Doğal tümdengelim !!{derivation} ağacının
yapraklarındaki formüllere varsayım denir ve bazı çıkarım kuralları bunların
belirli örnekleri için “!!{discharge}” işlemini uygulayabilir. Örneğin, aralarında $!A$'nın da
bulunduğu bazı varsayımlardan $!B$ için !!a{derivation} varsa
$\Intro{\lif}$ kuralı,
$!A \lif !B$ sonucunu çıkarmamıza ve $!A$ biçimindeki seçili varsayımları
kapatmamıza izin verir. (Hangi varsayımların hangi çıkarımda kapatıldığını
izlemek için çıkarımı ve onun kapattığı varsayımları bir sayıyla etiketleriz.)
!!{undischarged} kalan varsayımlar, !!{derivation} sonunda birlikte sonucun
doğruluğu için yeterlidir; böylece !!a{derivation}, !!{undischarged}
varsayımlarının sonucunu mantıksal olarak gerektirdiğini ortaya koyar.

Doğal tümdengelime dayalı $\Gamma \Proves !A$ bağıntısı, !!a{derivation}
içinde $!A$ ağaçtaki son !!{sentence} ise ve !!{undischarged} kalan her yaprak
$\Gamma$'da bulunuyorsa geçerlidir. !!a{derivation} içinde $!A$ son
!!{sentence} ise ve bütün varsayımlar !!{discharged} ise
$!A$, doğal tümdengelimde bir teoremdir. Örneğin aşağıdaki !!a{derivation},
$\Proves (!A \land !B) \lif !A$ olduğunu gösterir:
\begin{prooftree}
  \AxiomC{$\Discharge{!A \land !B}{1}$}
  \RightLabel{\Elim{\land}}
  \UnaryInfC{$!A$}
  \DischargeRule{\Intro{\lif}}{1}
  \UnaryInfC{$(!A \land !B) \lif !A$}
\end{prooftree}
$1$ etiketi, $!A \land !B$ varsayımının $\Intro{\lif}$ çıkarımında
!!{discharged} olduğunu gösterir.

Doğal tümdengelimde bir $\Gamma$ kümesi, ancak ve ancak $\Gamma \Proves
\lfalse$ ise tutarsızdır. \FalseInt{} kuralı sayesinde tutarsız bir kümeden
her bir !!{sentence} için onu !!{derive} mümkündür.

Doğal tümdengelim sistemleri 1930'larda Gerhard Gentzen ve Stanis\l{}aw
Ja\'skowski tarafından geliştirildi; daha sonra Dag Prawitz ve Frederic Fitch
tarafından ileri götürüldü. Çıkarımları doğal ispat yöntemlerini yansıttığı
için filozoflarca tercih edilir. Fitch'in geliştirdiği sürümler giriş düzeyi
mantık ders kitaplarında sıkça kullanılır. Mantık felsefesinde doğal
tümdengelimin kurallarının bazen mantıksal işleçlere anlamlarını verdiği kabul
edilmiştir (“ispat kuramsal semantik”).

\end{document}
